\documentclass[10pt,a4paper]{amsart}
\usepackage[margin=19mm]{geometry}
\usepackage{amsmath,amssymb,mathtools}
\usepackage[scr=rsfs,cal=euler]{mathalfa}
\usepackage{xfrac}
\usepackage[unicode,hidelinks,bookmarks=false]{hyperref}
\newcommand{\ie}{i.e.\ }
\newcommand{\cf}{cf.\ }
\newcommand{\resp}{respectively\ }
\newcommand{\Subsection}{\subsection}
\providecommand{\nobreakdash}{\nobreak\hbox{-}}
\setlength{\emergencystretch}{2em}
\multlinegap0pt
\usepackage{bm}
\usepackage{bbm}
\usepackage{dsfont}
\usepackage{xspace}
\usepackage{cancel}
\usepackage{mathtools}
\usepackage{amsthm}
\makeatletter
\@ifundefined{theo}{\newtheorem{theo}{Theo}}{}
\@ifundefined{defi}{\newtheorem{defi}[theo]{Definition}}{}
\@ifundefined{prop}{\newtheorem{prop}[theo]{Proposition}}{}
\@ifundefined{rem}{\newtheorem{rem}[theo]{Remark}}{}
\makeatother
\title[Mean field games with terminal state constraints]{Mean field games with terminal state constraints}
\author{Luciano Campi, Luca Di Persio, Viktorya Vardanyan}
\date{Source: arXiv:2606.11493v2 (CC BY 4.0)}
\begin{document}
\maketitle

\noindent\emph{Source and licence.} Mean field games with terminal state constraints. Version arXiv:2606.11493v2 (CC BY 4.0), distributed under the Creative Commons Attribution 4.0 International licence, as recorded in the arXiv licence metadata for this exact version and shown on the abstract page https://arxiv.org/abs/2606.11493v2. Full rights notice: licenses/arxiv-2606.11493-CC-BY-4.0.txt.\par\smallskip

\noindent\textit{Editorial scope.} Counted unit: the Prop whose source label is \texttt{prop:application1} in the article (source0.tex lines 1003--1016, source label \texttt{prop:application1}), delivered together with the article's own complete original proof of it (source0.tex lines 1018--1184, one \texttt{proof} environment of 167 lines). No mathematics is written, repaired or re-derived: statement and proof are the article's own text line for line. Required context retained uncounted: def-MFG-constraint (lines 244--260); def-auxMFG (lines 350--364); auxoptimal (lines 353--356); equivalence (lines 367--370); verification (lines 493--513); uniquexi (lines 645--648); uniq-appl (lines 737--743); Lasry (lines 747--765); Lasry-appl (lines 768--774); FBSDE\_Cont CN l (lines 812--824); EXUN (lines 827--830); remExun (lines 958--961); sde oi (lines 993--996). Every definition, display and label of those lines is retained unchanged. Omitted from this package and not redistributed: 1--243, 261--349, 357--366, 371--492, 514--644, 649--736, 744--746, 766--767, 775--811, 825--826, 831--957, 962--992, 997--1002, 1017--1017, 1185--2230. Nothing inside a retained excerpt is omitted or altered: every retained body is delivered exactly as the article's lines are written, so no omission receipt of any kind accompanies this package. Citations: the retained excerpts carry no citation command at all, so no bibliography entry of the article is transcribed and no reference is redistributed. Reference adaptation: none was needed and none was made; every reference inside the delivered unit resolves inside it, so no unresolved reference or citation is produced. Printed slips kept literal: no printed word, symbol, hypothesis or number of the counted statement or of its proof was altered. Layout and class: the article's own class is \texttt{article} with packages including babel, inputenc, setspace, geometry, epsfig, titlesec, sectsty, enumerate, amsmath, amsthm, amsfonts, amssymb, calrsfs, cases; the wrapper is the lane's pinned \texttt{amsart} editorial wrapper at 10pt on A4, which restates the theorem-like environments the retained text needs and adds the article's own macro definitions that the retained text needs (none), with extra wrapper packages (bm, bbm, dsfont, xspace, cancel, mathtools). The delivered numbering is the unit's own. Assets: no retained span contains a figure environment, a graphics-inclusion command, a PDF-page inclusion, a table or a TikZ picture; no image file is copied into this package, the package declares no asset dependency and no image file is redistributed with it. Licence: version arXiv:2606.11493v2 of this article is distributed under the Creative Commons Attribution 4.0 International licence, as recorded in the arXiv licence metadata for this exact version and shown on the abstract page https://arxiv.org/abs/2606.11493v2. A full rights notice is delivered at licenses/arxiv-2606.11493-CC-BY-4.0.txt. Characters: every retained excerpt is pure ASCII, so the delivered engine, XeLaTeX with its default Latin Modern text font, reports no missing character.
\begin{defi} \label{def-MFG-constraint} A MFG solution with terminal state constraint is any tuple $(\hat \alpha, \hat\beta, \hat \mu)$ of processes such that
\begin{enumerate}
\item $(\hat\alpha, \hat\beta) \in \mathcal A$ and  $\hat \mu$ is a $(\mathcal F_t ^0)$-progressively measurable process with values in $\mathcal{P}_2(\mathbb R^n \times \mathbb R^{n\times d} \times \mathbb R^{n\times d})$, such that $\mathbb E \left[\int_0^T M_2(\hat \mu_t)^2 dt\right]<\infty$;
\item (Optimality condition) Given the flow $\hat \mu$, the pair $(\hat \alpha, \hat \beta)$ solves the following optimal control problem
\begin{equation}
\label{forward CN}
     \inf_{(\alpha, \beta)\in \mathcal A,  X_T \in \mathcal D} \mathbb{E}\left[\int _0^T \bar{\ell}(t,X_t,  \alpha_t, \beta_t, \hat\mu _t)dt+\phi(X_T,\hat\mu^X_T)\right],
\end{equation}
where
\begin{equation}
\begin{aligned}
    dX_t=b(t,X_t,  \alpha_t, \beta_t, \hat\mu_t)dt+\sigma(t,X_t, \alpha_t, \hat\mu_t)dW_t
    +\sigma_0(t,X_t, \beta_t, \hat\mu_t)dW_t^0, \quad  X_0=a;
 \end{aligned}   
\end{equation}
\item (Consistency condition)  For all $t\in [0,T]$, we have \[\hat\mu_t=\mathcal{L}(\hat X_t, \hat\alpha_t, \hat\beta_t|\mathcal{F}_t^0),\] where $\hat X$ denotes the state controlled by $(\hat\alpha, \hat\beta)$.    \end{enumerate}
\end{defi}

\begin{defi} \label{def-auxMFG} A solution to the \emph{auxiliary} MFG with terminal state constraint is any pair $(\hat  \xi, \hat \mu)$, where $\hat \xi \in L^2 _a (\mathcal D)$, $\hat \mu$ is an $(\mathcal F_t ^0)$-adapted process with values in $\mathcal{P}_2(\mathbb{R}^n\times \mathbb{R}^{n\times d} \times \mathbb{R}^{n\times d})$, such that $\mathbb E \left[\int_0^T M_2(\hat \mu_t)^2 dt\right]<\infty$, and the following two properties are satisfied:
\begin{enumerate}
    \item (Optimality condition) Given $\hat \mu$, the random variable $\hat \xi$ solves the following optimal control problem
\begin{equation}
\label{auxoptimal}
     \inf_{\xi\in L^2 _a (\mathcal D)}\mathbb E \left[\int_0^T \ell(t, X_t, q_t, z_t, \hat \mu_t)dt+\phi(\xi, \hat\mu^X_T)\right],
\end{equation}
where $(X,q,z) \in \mathcal S^2_n \times  \mathcal H^2_{n \times d}\times \mathcal H^2_{n\times d}$ solves
\begin{equation}
-dX_t=f(t, X_t, q_t, z_t, \hat \mu_t)dt -q_tdW_t-z_tdW_t^0, \quad X_T=\xi;
\end{equation}
\item (Consistency condition)  For all $t\in [0,T]$, we have \[\hat \mu_t=\mathcal{L}(\hat X_t, \hat\alpha_t, \hat\beta_t |\mathcal{F}_t^0), \quad t\in [0,T], \] where $\hat X$ denotes the state with terminal condition $\hat X_T  = \hat \xi $ and
\begin{equation} (\hat \alpha_t , \hat \beta_t ) = (\tilde \sigma, \tilde \sigma_0)(t, \hat X_t , \hat q_t , \hat z_t , \hat \mu_t ). \label{eq:xi-to-alpha-beta}\end{equation}
\end{enumerate}
\end{defi}

\begin{prop} 
\label{equivalence}
Under assumptions (A1)-(A5), a triple $(\hat \alpha, \hat \beta, \hat \mu)$ is a solution of the MFG with terminal state constraint  (Definition \ref{def-MFG-constraint}) if and only if the pair $(\hat \xi, \hat \mu)$ is a solution of the auxiliary MFG as in Definition \ref{def-auxMFG}.
\end{prop}

\begin{theo} 
\label{verification}
Suppose that $H(t,\cdot, \cdot, \cdot, \cdot, \mu, y, h_0)$ and $\phi(\cdot, \mu)$ are convex in $(x,q, z)$ and $x$ respectively. 
Then any pair $(\hat \xi, \hat \mu )$ is a solution for the auxiliary MFG with terminal state constraint as in Definition \ref{def-auxMFG} if the tuple $( \hat X,   \hat q,  \hat z,  \hat \mu,   \hat Y,  \hat {h}_0, \hat {h}_1, \hat {\xi})$, where $( \hat X,   \hat q,  \hat z,  \hat Y) \in \mathcal S^2 _n \times \mathcal H^2 _{n\times d} \times \mathcal H^2 _{n\times d} \times \mathcal S^2 _n$, $\hat \mu$ is an $(\mathcal F_t ^0)$-progressively measurable process such that $\mathbb E \left[\int_0^T M_2(\hat \mu_t)^2 dt\right]<\infty$, $( \hat h_0,  \hat h_1) \in \mathbb R \times \mathbb R^n$, and $\hat \xi \in L^2(\mathcal D)$, solves the FBSDE system
\begin{equation}
\label{FBSDE MFG CN}
\begin{cases}
    -dX_t=f(t,X_t, q_t, z_t, \mu_t)dt -q_tdW_t-z_tdW_t^0\\
      X_0 =a, \quad X_T=\xi\\
    dY_t=[f_x(t)Y_t+h_0 \ell_x(t)]dt+[f _q(t)Y_t+h_0 \ell_q(t)]dW_t+[f_z(t)Y_t+h_0 \ell_z(t)]dW_t^0\\
        Y_0=h_1\\
        h_0 > 0\\
    \big( Y_T+h_0\phi_x(\xi,\mu_T^X)\big) \big(\xi^\prime -\xi \big) \geq 0 \ a.s., \ \forall \xi^\prime \in L^2(\mathcal{D})
\end{cases}
\end{equation}
with
\[\hat{\mu}_t=\mathcal{L}(\hat X_t, \hat \alpha_t , \hat {\beta}_t \mid \mathcal F_t ^0),   \quad  t\in [0,T],\]
where 
$ (\hat \alpha_t , \hat \beta_t ) = (\tilde \sigma, \tilde \sigma_0)(t, \hat X_t , \hat q_t , \hat z_t , \hat \mu_t ).$
The reverse also holds true with the condition $h_0 >0$ in the FBSDE system replaced with the weaker $h_0 \geq 0$.
\end{theo}

\begin{prop}
\label{uniquexi}
Let assumptions (C1)-(C4) hold. Then the optimal control problem \eqref{auxoptimal} admits at most one minimizer in $L^2_a(\mathcal D)$.
\end{prop}

\begin{rem}\label{uniq-appl}
In the applications of Section 4,  the terminal cost functional
\[
\phi(x,\bar{\mu}) = -x + \frac{\gamma}{2}(x - \theta\bar{\mu})^2
\]
is uniformly $\lambda$-convex in $x \in \mathbb{R}$, uniformly in $\bar{\mu}$, with any $\lambda \in (0, \gamma / 2]$. 
\end{rem}

\begin{prop}
\label{Lasry}
 Let assumptions \textnormal{(C1)--(C3), (C5)} hold,  and  suppose that for any given $(\mathcal{F}_t^0)$-adapted process $\hat \mu^X =(\hat \mu^X _t)_{0\leq t\leq T}$ with values in  $\mathcal{P}_2(\mathbb{R}^n) $, the optimal control problem 
 \begin{equation} \label{eq:opt_problem}
\inf_{\xi \in L^2_a(\mathcal{D})} \, \mathbb{E} \Bigg[  \int_0^T \ell(t, X_t, q_t, z_t ,  \hat \mu_t^X) \, dt + \phi(\xi, \hat \mu^X_T) \Bigg],
\end{equation}
subject to
\begin{equation}
\begin{cases}
    -dX_t =  \, f(t, X_t, q_t, z_t) \, dt - q_t \, dW_t-z_tdW_t^0\\  
X_0 =  \, a, \quad X_T = \xi,
\end{cases}
\end{equation}
admits a unique minimizer \( \hat \xi\in L^2_a(\mathcal{D})\). Let $( \hat X, \hat q, \hat z )$ denote the corresponding optimal state. Then, there exists at most one flow $\hat \mu^X=(\hat \mu_t^X)_{0\leq t\leq T}$ such that
\begin{equation}
\label{cons}
  \hat \mu_t^X=\mathcal{L}(\hat X_t|\mathcal{F}_t^0), \quad t \in [0, T]. 
\end{equation}
\end{prop}

\begin{rem}\label{Lasry-appl} 
We observe that in the applications of Section 4, the terminal cost functional
\[
\phi(x,\mu) = -x + \frac{\gamma}{2}\left(x - \theta \int_{\mathbb R}x\,d\mu(x)\right)^2
\]
is not monotone in the sense of Lasry-Lions. Consequently, Proposition \ref{Lasry} is not applicable. Hence, in order to get uniqueness we will rely on a contraction argument instead.
\end{rem}

\begin{equation}
\label{FBSDE_Cont CN l}
    \begin{cases}
    dX_t=-[b_0(t)X_t+q_tc_0(t)^\top +z_t\tilde{c}_0(t)^\top]dt +q_tdW_t+z_tdW_t^0\\
      X_0 =a, \quad X_T=\xi\\
  
    dY_t=b_0(t)Y_tdt+c_0(t)Y_tdW_t
    +\tilde{c}_0(t)Y_tdW_t^0\\
   Y_0=h_1\\

    \big( Y_T+ \phi_x(\xi, m_T)\big)\big( \xi'-\xi \big) \geq 0 \ a.s., \ \forall \xi' \in L^2([c, \infty)).        
   \end{cases}
\end{equation}

\begin{theo}
\label{EXUN} 
Let $a > c \ e^{\int_0^T b_0(t)\, dt}$.    Under Assumptions  (A1)-(A7), the FBSDE system \eqref{FBSDE_Cont CN l} has a unique  solution.
\end{theo}

\begin{rem}
\label{remExun}
Similar results as in Theorem \ref{EXUN} (and its proof) hold also in the two ``extreme'' cases with no common noise and with only common noise with the corresponding modifications. They will be used in the applications of the next section.
\end{rem}

\begin{equation}
\label{sde oi}
    dX_t = \alpha_t X_t (b dt + \sigma dW_t ), \quad X_0=a, \quad X_T \geq c,
\end{equation}

\begin{prop}\label{prop:application1}
Let $(\hat h_1, \hat m_T) \in \mathbb R \times \mathbb R$ be the unique solution of the coupled system 
\begin{equation} \label{eq:h1}
\begin{cases}
 m_T = c + \mathbb{E} \left[ \left( \frac{1}{\gamma}-c+\theta m_T - \frac{h_1}{\gamma} Y_T^{(1)} \right)^{+} \right]\\
    G(h_1):= c -a+ \mathbb{E} \left[ Y_T^{(1)} \left( \frac{1}{\gamma}-c+\theta m_T - \frac{h_1}{\gamma} Y_T^{(1)} \right)^{+} \right]=0,\\
\end{cases}
\end{equation} 
 and let
\[ \hat g(t,y) := \mathbb{E}\!\left[ yY^{(1)}_{t,T} \left(\frac{1}{\gamma}-c+\theta \hat{m}_T- \frac{\hat h_1 }{\gamma}yY^{(1)}_{t,T} \right)^+ \right], \quad (t,y)\in [0,T]\times \mathbb R^+.\]
 The unique MFG solution is given by the couple $(\hat \alpha, \hat m_T)$ where
$$\hat\alpha_t=-\frac{\rho}{\sigma}
\dfrac{Y^{(1)}_t\partial_{y}\hat g(t, Y^{(1)}_t) - \hat g(t,Y^{(1)}_t)}{cY^{(1)}_t+\hat g(t, Y^{(1)}_t)}, \quad t\in [0,T).$$
\end{prop}

\begin{proof}
First, we reformulate the control problem in terms of the amount $\pi_t$ invested
in the risky asset at time $t$, i.e., by letting
\[
\pi_t:=\alpha_tX_t, \quad t\in [0,T],
\]
the wealth dynamics \eqref{sde oi} becomes
\[
dX_t=\pi_t(b\,dt+\sigma\,dW_t),\qquad X_0=a.
\]
In this reformulated problem the diffusion coefficient \(\pi \sigma\) satisfies assumption (A4), so that we can use the approach developed in the previous sections and set
\[
q_t:=\sigma\pi_t,
\qquad \rho:=\frac b\sigma,
\]
yielding
\[
dX_t=q_t(\rho\,dt+dW_t),\qquad X_0=a.
\]
After solving the reformulated problem, the original control will be recovered via
\[
\alpha_t=\frac{\pi_t}{X_t}=\frac{q_t}{\sigma X_t},
\]
whenever \(X_t>0\), which we will verify to be true at the equilibrium. Using Theorem \ref{verification}, we are reduced to solving the following FBSDE system:
\begin{equation}
\label{FBSDE general 1 oi}
\begin{cases}
    dX_t=  q_t (\rho dt+ dW_t)\\
    X_0=a, \quad X_T=\xi\\
     dY_t=-\rho Y_t dW_t\\
    Y_0=h_1\\
        
       
        \big(Y_T-1+\gamma(\xi -\theta m_T)\big)\big( \xi'-\xi \big) \geq 0 \ \text{a.s.}, \ \forall \xi' \in L^2([c,+\infty)) .
\end{cases}
\end{equation}
Applying the same computations as in the first part of Theorem \ref{EXUN} (see also Remark \ref{remExun}) with the parameters
\[
c_0(t) = -\rho, \quad b_0(t) = 0, \quad \tilde{c}_0(t) = 0, \quad k_0 = -1, \quad k_1 = \gamma, \quad k_2 =-\gamma\theta,
\]
we have that the terminal condition of system \eqref{FBSDE general 1 oi} takes the form 
\begin{equation}
\label{terminal_X}
\xi=c + \left( \frac{1}{\gamma} - c + \theta m_T - \frac{h_1}{\gamma} Y_T^{(1)} \right)^{+}.
\end{equation}
By the same arguments as in the proof of Theorem \ref{EXUN}, we deduce that there exists a unique  pair $(\hat h_1, \hat m_T)$  which solves the following  coupled system
\begin{align*}
m_T &= c + \mathbb{E} \left[ \left( \frac{1}{\gamma} - c + \theta m_T - \frac{h_1}{\gamma} Y_T^{(1)} \right)^{+} \right], \\
a &= c + \mathbb{E} \left[ Y_T^{(1)} \left( \frac{1}{\gamma} - c + \theta m_T - \frac{h_1}{\gamma} Y_T^{(1)} \right)^{+} \right].
\end{align*}
Since \(X Y^{(1)}\) is a martingale, we have
\[
X_t = \frac{1}{Y_t^{(1)}}
\mathbb{E}\left[ X_T Y_T^{(1)} \mid \mathcal{F}_t \right].
\]
Using the terminal condition
\[X_T=c+\left(\frac{1}{\gamma} - c + \theta m_T-\frac{h_1}{\gamma} Y_T^{(1)}\right)^+,
\]
we obtain
\[
X_t=c+\frac{1}{Y_t^{(1)}}\mathbb{E}\left[Y_T^{(1)}
\left(\frac{1}{\gamma} - c + \theta m_T-\frac{h_1}{\gamma} Y_T^{(1)}\right)^+\;\middle|\;\mathcal{F}_t\right].\]
We notice that, since $Y^{(1)}_T = Y^{(1)}_t Y^{(1)}_{t,T}$, where $Y^{(1)}_t$ and $Y^{(1)}_{t,T}$ are independent, we can express $X$ as follows
\[
X_t = c + \frac{1}{Y^{(1)}_t} \hat g(t, Y^{(1)}_t),
\]
where
\[
\hat g(t,y) := \mathbb{E}\!\left[ yY^{(1)}_{t,T} \left(\frac{1}{\gamma} - c + \theta \hat m_T- \frac{\hat h_1 }{\gamma}yY^{(1)}_{t,T} \right)^+ \right].\]
Observe that $Y^{(1)}_{t,T}$ follows a lognormal distribution with a smooth density on $(0,\infty)$. Hence, $\hat g(t,\cdot)$ is a lognormal convolution of a payoff with at most quadratic growth, so $\hat g \in C^{1,2}([0,T) \times (0, \infty))$. Since $Y^{(1)} > 0$, we can safely apply It\^o's formula to $X_t = c+  \hat g(t,Y^{(1)}_t)/Y^{(1)}_t$, $t<T$, yielding
\[
dX_t = \left(\partial_t \kappa (t,Y^{(1)}_t) + \frac12 \rho^2 (Y_t ^{(1)})^2 \, \partial_{yy} \kappa (t,Y^{(1)}_t) \right)dt - \rho Y^{(1)}_t \, \partial_y \kappa (t,Y^{(1)}_t) dW_t, \quad t<T,
\]
where $\kappa (t,y) :=y^{-1}\hat g(t,y)$. Therefore, we can identify the process $q$ as
\[
q_t=-\rho Y_t^{(1)} \partial_y \kappa (t,Y^{(1)}_t)=-\rho
\left(
\partial_{y}\hat g(t,Y^{(1)}_t) - \frac{1}{Y^{(1)}_t} \hat g(t,Y^{(1)}_t)\right),\quad   t<T.
\]
Finally, since
\[
X_t=c+\frac{\hat g(t,Y_t^{(1)})}{Y_t^{(1)}}>0,
\]
\(\hat\alpha_t\) is well defined, and recalling that
\(q_t=\alpha_tX_t\sigma\), we obtain the equilibrium strategy
\begin{equation}
\label{opt-alpha}
\hat \alpha_t= -\frac{\rho}{\sigma}
\dfrac{Y^{(1)}_t\partial_{y}\hat g(t,Y^{(1)}_t) - \hat g(t,Y^{(1)}_t)}{cY^{(1)}_t+\hat g(t,Y^{(1)}_t)}, \quad t\in [0,T).
\end{equation}
It remains to prove uniqueness. For notational convenience, we write $m$ instead of $m_T$.	
Fix $m\in \mathbb R$. By Proposition \ref{uniquexi} and Remark \ref{uniq-appl}, $\xi^m\in L^2_a([c, +\infty))$ given by \eqref{terminal_X} for the fixed $m$, is  the unique minimizer of 
\[
J_m(\xi) := \mathbb E\left[-\xi+\frac{\gamma}{2}(\xi-\theta m)^2\right]. 
\]
As observed in  Remark~\ref{Lasry-appl}, the terminal cost functional is not monotone in the sense of Lasry-Lions. Therefore, we establish uniqueness via the following contraction argument.
Define
\[
\Gamma(m):=\mathbb E[\xi^m].
\]
We show that \(\Gamma\) is a contraction.
Let \(m_1,m_2\in\mathbb R\), and set
$\xi_1:=\xi^{m_1}, \xi_2:=\xi^{m_2}$.
Since \(L^2_a([c, +\infty))\) is convex, for every \(\varepsilon\in(0,1)\),
\[
\xi_\varepsilon:=(1-\varepsilon)\xi_1+\varepsilon\xi_2\in L^2_a([c, +\infty)).
\]
Moreover, since \(\xi_1\) minimizes \(J_{m_1}\) over \(L^2_a([c, +\infty))\),
\[
J_{m_1}(\xi_1)\le J_{m_1}(\xi_\varepsilon).
\]
On the other hand, by the \(\lambda\)-convexity of \(J_{m_1}\),
\[
J_{m_1}(\xi_\varepsilon)\le(1-\varepsilon)J_{m_1}(\xi_1)+\varepsilon J_{m_1}(\xi_2)-\lambda\varepsilon(1-\varepsilon)\mathbb E\left[|\xi_2-\xi_1|^2 \right].
\]
Combining these last two inequalities gives
\[
J_{m_1}(\xi_1)\le(1-\varepsilon)J_{m_1}(\xi_1)+\varepsilon J_{m_1}(\xi_2)-\lambda\varepsilon(1-\varepsilon)\mathbb E\left[|\xi_2-\xi_1|^2 \right].
\]
Rearranging and dividing by \(\varepsilon>0\), we obtain
\[
J_{m_1}(\xi_2)-J_{m_1}(\xi_1)\ge \lambda(1-\varepsilon)\mathbb E\left[|\xi_2-\xi_1|^2 \right].
\]
Repeating the same argument with \(m_2\), minimizer \(\xi_2\), and competitor \(\xi_1\), yields
\[
J_{m_2}(\xi_1)-J_{m_2}(\xi_2)\ge\lambda (1-\varepsilon)\mathbb E\left[|\xi_2-\xi_1|^2 \right].
\]
Adding the last two inequalities gives
\begin{equation}\label{ineq-m1m2}
2\lambda (1-\varepsilon)\mathbb E \left[|\xi_2-\xi_1|^2 \right]\le
J_{m_1}(\xi_2)-J_{m_1}(\xi_1)+J_{m_2}(\xi_1)-J_{m_2}(\xi_2).
\end{equation}
Since
\[
J_m(\xi)=\mathbb E\left[-\xi+\frac{\gamma}{2}(\xi-\theta m)^2\right],
\]
the right-hand side in \eqref{ineq-m1m2} satisfies
\[
\begin{aligned}
&J_{m_1}(\xi_2)-J_{m_1}(\xi_1)+J_{m_2}(\xi_1)-J_{m_2}(\xi_2) \\
&=\frac{\gamma}{2}\mathbb E\Big[(\xi_2-\theta m_1)^2-(\xi_1-\theta m_1)^2+(\xi_1-\theta m_2)^2-(\xi_2-\theta m_2)^2\Big]\\
&=\gamma\theta(m_1-m_2)\mathbb E[\xi_1-\xi_2].
\end{aligned}
\]
Therefore,
\begin{align*}
2\lambda (1-\varepsilon)\mathbb E \left[|\xi_2-\xi_1|^2 \right] & \le \gamma\theta(m_1-m_2)\mathbb E[\xi_1-\xi_2]\\
& \le \gamma\theta |m_1-m_2|\left|\mathbb E[\xi_1-\xi_2]\right|.
\end{align*}
By Jensen's inequality,
$\left|\mathbb E[\xi_1-\xi_2]\right|^2 \le \mathbb E\left[|\xi_1-\xi_2|^2\right]$.
Hence, if \(\mathbb E[\xi_1-\xi_2]\neq0\), we obtain
\[
\left|\mathbb E[\xi_1-\xi_2]\right| \le\frac{\gamma\theta}{2\lambda(1-\varepsilon)}|m_1-m_2|.
\]
If \(\mathbb E[\xi_1-\xi_2]=0\), the same inequality trivially holds too. Thus
\[
|\Gamma(m_1)-\Gamma(m_2)|\le\frac{\gamma\theta}{2\lambda(1-\varepsilon)}|m_1-m_2|.
\]
By Remark \ref{uniq-appl}  for the present quadratic cost we may take \(\lambda=\gamma/2\). Therefore,
\[
|\Gamma(m_1)-\Gamma(m_2)| \le \frac{\theta}{1-\varepsilon} |m_1-m_2|.
\]
Since \(\theta\in[0,1)\), we may choose \(\varepsilon>0\) sufficiently small such that $\frac{\theta}{1-\varepsilon}<1$, yielding that \(\Gamma\) is a contraction on \(\mathbb R\). Therefore the consistency condition
$m=\mathbb E[\xi^m]$
has at most one solution. Finally, by Proposition \ref{equivalence} we have  uniqueness of the MFG equilibrium $(\hat \alpha, \hat m_T)$.
\end{proof}

\end{document}
